\section{The renewal passage at unchanged memory}\label{sec:proof}
\begin{lemma}[Finite-response continuity]\label{lem:tensor}
Equip each simplex-valued report row with its weak-star $L^\infty(\sigma_a)$ topology, each action row with its Euclidean topology, and $D^M$ with its compact topology. The resulting space $\Ga$, modulo the common reference-null sets, is compact. For fixed $\theta,T$, the maps
\[
 \gam\longmapsto\E_\theta^\gam(\tau\wedge T),\qquad
 \gam\longmapsto\E_\theta^\gam\sum_{t\le\tau\wedge T}\ell_t
\]
are continuous.
\end{lemma}
\begin{proof}
Positivity and the row-sum identity are weak-star closed constraints in finitely many unit balls. Banach--Alaoglu and the finite product theorem give compactness. On a standard Borel probability space, a completed-measurable representative has a Borel version. Complete all finitely many rows on one union of reference-null exceptional sets by the point mass at $1$. Product domination makes this completion common to every model and finite history.

For clarity, the continuity argument applies to nets. If $u_{j,\lambda}\to u_j$ weak-star and $|u_{j,\lambda}|\le1$, then for any $f\in L^1(\bigotimes_{j=1}^t\sigma_j)$,
\[
 \int f(y)\prod_{j=1}^tu_{j,\lambda}(y_j)d\sigma(y)
 \longrightarrow\int f(y)\prod_{j=1}^tu_j(y_j)d\sigma(y).
\]
It holds first for rectangle indicators, where the integral factors, and then for their finite linear combinations. These are $L^1$-dense, and the replacement error is uniform in the net. Several factors may be copies of one changing row, because their report coordinates are distinct.

For each $t$, form the payoff submeasure
\[
 \Lambda_{\theta,t}^{\mathbf a,d}(B)
 =\E_\theta^{\operatorname{do}(\mathbf a)}
       [1_{\{\tau\ge t,\,Y_{1:t}\in B\}}\ell(d,Z_t)].
\]
It is dominated by the report measure in \eqref{eq:product}. Its density is $L^1$-continuous in $d$ by bounded convergence. Causal expansion gives, with $i_0=1$,
\begin{equation}\label{eq:expansion}
 \E_\theta^\gam[1_{\{\tau\ge t\}}\ell_t]
 =\sum_{\mathbf a,\mathbf i}\int
 f_{\theta,t}^{\mathbf a,d_{i_t}}(y)
 \prod_{s\in\mathcal D_t}\alpha_{i_{s-1}}(a_s)
 k_{i_{s-1},a_s}(i_s\mid y_s)\,d\sigma_{\mathbf a}(y).
\end{equation}
Here $\mathcal D_t$ consists of the ordinary data positions; forced preparation positions have their prescribed factor. The action probabilities are in the product. The expansion sums interventions; it does not condition on an event of possibly zero probability. Tensor continuity proves continuity of each finite sum. Using the survival submeasure in place of the payoff and summing over $t\le T$ proves the first assertion too. This is the reused-row step, rather than continuity of pointwise multiplication in one $L^\infty$ space \cite{Discounted}.
\end{proof}

\begin{remark}[Why the product hypothesis matters]\label{rem:diagonal}
For uniform $Y\in[0,1]$, let two reports both equal $Y$, and let $r_n$ be successive dyadic $\{0,1\}$-valued oscillations. Then $r_n\to1/2$ weak-star, but $\E r_n(Y)^2=1/2$, not $1/4$. The report pair is diagonal and singular to product Lebesgue measure. Independent-position static reduction does not justify the missing limit. Conversely, after product domination the diagonal constraint that the same row is reused is closed in the product row topology, and Lemma~\ref{lem:tensor} supplies its continuity. The subsequent compact minimization is standard; no new abstract compactness principle is claimed.
\end{remark}

\begin{proof}[Proof of Theorem~\ref{thm:renewal}]
The truncation errors for $r_\theta$ and $\mu_\theta$ are at most $a(T)$. They are therefore continuous by Lemma~\ref{lem:tensor} and uniform convergence. Since $\mu_\theta\ge1$, so is $g_\theta$. The supremum of these continuous functions is lower semicontinuous on one compact space, and attains its minimum. For a level $u$, the sets $\{\gam:g_\theta(\gam)\le u\}$ are closed. If every finite subfamily admits a point in their intersection, compactness gives a point in the full intersection. Applying this to levels above the supremum of finite-subfamily values proves \eqref{eq:finite-model}. The expansion proves response invariance.

For fixed $\theta,\gam$, cycles $(\tau_j,R_j)$ are iid, conditional on this fixed model. Write $T_k=\sum_{j=1}^k\tau_j$ and $K_N=\min\{k:T_k>N\}$. Since $\tau_j\ge1$, $K_N\le N+1$. The centered cycle increments $X_j=R_j-g_\theta\tau_j$ have mean zero. The event $\{K_N\ge j\}=\{T_{j-1}\le N\}$ depends only on preceding cycles; consequently
\[
 \E\sum_{j=1}^{K_N}X_j=0.
\]
Removing from the last cycle its reward and time after call $N$ changes this centered sum by a random variable whose absolute value is at most $T_{K_N}-N$. Thus
\begin{equation}\label{eq:overshoot}
 |\E C_N-Ng_\theta|\le\E(T_{K_N}-N)
 \le\sum_{s=0}^N\E(\tau-s)_+.
\end{equation}
For the last inequality, sum over the last renewal epoch before $N$. A given integer is a renewal epoch with probability at most one, because increments are positive integers. This proves the first bound in \eqref{eq:uniform}. This elementary bound is consistent with the classical overshoot theory \cite{Lorden}; we do not require a finite second moment for the tail-envelope form.

The iid strong laws give $T_k/k\to\mu_\theta$ and $\sum_{j\le k}R_j/k\to r_\theta$. Also $\tau_k/k\to0$ almost surely, by the tail-sum criterion and Borel--Cantelli. The residual cycle contributes at most its length, so $C_N/N\to g_\theta$ almost surely; boundedness gives convergence in mean. These are fixed-model, fixed-program assertions.

Abel summation gives the exact identity
\begin{equation}\label{eq:abel}
 J_{\theta,\beta}(\gam)
 =(1-\beta)^2\sum_{N\ge1}N\beta^{N-1}J_{\theta,N}(\gam).
\end{equation}
The coefficients sum to one. Equation~\eqref{eq:uniform} and the definition of $c_\beta$ follow. Since $b_N\to0$ and it is bounded, $c_\beta\to0$. Taking suprema and infima preserves a common uniform error; it does not change the program class. Finite-prefix risks are continuous, and discounted risks are uniform limits of such risks with tail at most $\beta^T$, proving their attainment. If $\gam_j$ are asymptotically optimal for $N_j\to\infty$ or $\beta_j\uparrow1$, their average robust risks tend to $V_\av^\Theta(M)$. Lower semicontinuity at any convergent subnet proves the cluster-point assertion.

Finally, under \eqref{eq:moment}, the right side of \eqref{eq:overshoot} is bounded by
\[
 \E\sum_{s=0}^N(\tau-s)_+
 \le\E[\tau\min\{\tau,N+1\}]
 \le (N+1)^{1-p}\E\tau^{1+p}.
\]
For the discount bound, let $G$ have geometric distribution
$\PP(G=N)=(1-\beta)\beta^{N-1}$. Concavity gives
$\E(G+1)^{1-p}\le(\E G+1)^{1-p}\le[2/(1-\beta)]^{1-p}$.
Insert the preceding overshoot bound into \eqref{eq:abel} to obtain \eqref{eq:moment-rates}.
\end{proof}

\begin{proposition}[Sharp time orders and the need for a tail envelope]\label{prop:time-sharp}
The orders $N^{-p}$ and $(1-\beta)^p$ in \eqref{eq:moment-rates} are sharp uniformly over a bounded $1+p$ return-moment class, even with no control decisions. A uniform first-moment bound alone does not make average reward continuous under finite-response convergence.
\end{proposition}
\begin{proof}
A cycle has a zero-score preparation followed by either one zero-score data call, or $L$ unit-score calls. The latter event has probability $e=L^{-1-p}$. Then
\[
 \mu=2+e(L-1),\qquad g=\frac{eL}{2+e(L-1)},\qquad
 \E\tau^{1+p}\le 2^{1+p}+2^{1+p}eL^{1+p}.
\]
By call $N$ at most $N$ cycles can start, so $J_N\le Ne$. Taking $L=12N$ yields $J_N\le g/4$, since $\mu\le3$, and $g\ge L^{-p}/3$. For the discount, regeneration gives
\[
 J_\beta=\frac{e\beta(1-\beta^L)}
 {1-(1-e)\beta^2-e\beta^{L+1}}
 \le\frac{e}{1-\beta}.
\]
Choose an integer $L$ between $12/(1-\beta)$ and $13/(1-\beta)$ when $\beta$ is sufficiently close to one. Again $J_\beta\le g/4$ and $g$ has order $(1-\beta)^p$. These are uniform-class lower examples, not assertions that every fixed experiment attains either worst order.

For the last assertion instead take $e=1/L$. Finite response laws converge to the two-call zero-score cycle, while the mean lengths stay bounded by $3$ and $g\to1/3$. The limit cycle has average reward zero. The family fails uniform integrability exactly at its rare long cycles.
\end{proof}
